\chapter*{Lời nói đầu}
\addcontentsline{toc}{chapter}{Lời nói đầu}

Trong các hệ điều hành hiện đại, nhiều tiến trình và luồng thường cùng tồn tại, cùng chia sẻ bộ nhớ, thiết bị ngoại vi, tệp tin, hàng đợi công việc hoặc các tài nguyên dùng chung khác. Nếu việc truy cập các tài nguyên này không được điều phối đúng cách, hệ thống có thể rơi vào các lỗi khó phát hiện như tranh chấp dữ liệu, mất cập nhật, kẹt tiến trình, đói tài nguyên hoặc deadlock. Vì vậy, đồng bộ hóa tiến trình là một nội dung nền tảng trong học phần Nguyên lý Hệ điều hành.

Đề tài \textbf{``Các bài toán đồng bộ tài nguyên găng''} được thực hiện nhằm hệ thống hóa một số bài toán kinh điển thường dùng để mô phỏng các tình huống tranh chấp tài nguyên trong thực tế. Các bài toán như Roller Coaster, Room Party, River Crossing, Cigarette Smokers, Santa Claus, Sushi Bar, Dining Savages, One-Lane Bridge, Elevator, Ticket Selling, Parking Lot, Child Care, Bear and Honeybees, H2O, Multiplayer Matchmaking và Senate Bus đều đặt ra những ràng buộc khác nhau về số lượng tiến trình, thứ tự phục vụ, điều kiện ghép nhóm và quyền truy cập vào tài nguyên dùng chung.

Thông qua việc phân tích các bài toán trên, nhóm mong muốn làm rõ cách sử dụng các công cụ đồng bộ như mutex, semaphore, monitor và biến điều kiện để thiết kế lời giải đúng. Bên cạnh đó, báo cáo cũng chú trọng đến các tiêu chí đánh giá quan trọng như tính loại trừ lẫn nhau, tránh deadlock, tránh starvation và đảm bảo tiến trình được phối hợp theo đúng ràng buộc của bài toán.

Do giới hạn về thời gian và kinh nghiệm, báo cáo khó tránh khỏi thiếu sót trong cách trình bày, phân tích hoặc lựa chọn mô hình lời giải. Nhóm em rất mong nhận được nhận xét và góp ý của thầy để hoàn thiện hơn nội dung nghiên cứu.

Em xin chân thành cảm ơn!

\vspace{1.5cm}
\begin{flushright}
    \textit{Hà Nội, ngày ... tháng ... năm 2026} \\
    \vspace{0.2cm}
    \textbf{Nhóm sinh viên thực hiện} \\
\end{flushright}
